\documentclass[11pt,a4paper]{article}
\usepackage[utf8]{inputenc}
\usepackage[T1]{fontenc}
\usepackage[margin=2.1cm]{geometry}
\usepackage{amsmath,amssymb,amsthm}
\usepackage{booktabs,longtable,array}
\usepackage{enumitem}
\usepackage{xcolor}
\usepackage[colorlinks=true,linkcolor=blue!50!black,urlcolor=blue!50!black,citecolor=blue!50!black]{hyperref}
\setlength{\parskip}{3pt}
\newtheorem{theorem}{Théorème}[section]
\newtheorem{proposition}[theorem]{Proposition}
\newtheorem{lemma}[theorem]{Lemme}
\newtheorem{corollary}[theorem]{Corollaire}
\theoremstyle{remark}
\newtheorem{remark}[theorem]{Remarque}
\newcommand{\F}{\mathbb{F}}
\newcommand{\Z}{\mathbb{Z}}
\newcommand{\Q}{\mathbb{Q}}
\newcommand{\Fix}{\operatorname{Fix}}
\newcommand{\srg}{\operatorname{srg}}
\newcommand{\Aut}{\operatorname{Aut}}
\newcommand{\tr}{\operatorname{tr}}
\newcommand{\rg}{\operatorname{rang}}
\newcommand{\st}[1]{\textbf{[#1]}}
\newcommand{\code}[1]{\texttt{#1}}

\title{Le graphe de Conway $\srg(99,14,1,2)$, branche C3\\
\large Audit complet des travaux du 22 septembre 2026 : ce qui est démontré, certifié, calculé, et ce qu'il reste à faire}
\author{Travail et audit : Claude (Opus 5.5), pour D.~Djematene\\
\small Dossier de travail : \code{Desktop/c3\_recherche\_annexe} (hors de tout dépôt git)}
\date{22 septembre 2026}

\begin{document}
\maketitle

\begin{abstract}
On étudie l'hypothèse qu'un graphe fortement régulier $\Gamma=\srg(99,14,1,2)$ admette un automorphisme $g$ d'ordre $3$.
La littérature impose que $g$ soit sans point fixe. Ce document rassemble, avec leur statut exact, tous les résultats obtenus le 22 septembre 2026 :
(i) $g$ a exactement six orbites triangulaires (conséquence d'une identité d'Ishida, argument refait) ;
(ii) le graphe $F$ des adjacences entre ces six orbites est l'un de \textbf{26} graphes (classification certifiée par certificats exacts) ;
(iii) au niveau des appariements (« phases »), il reste \textbf{46} classes à symétrie près (certifié) ;
(iv) pour 35 de ces classes, le problème restant se ramène à une liste finie et explicite de problèmes « quotient », 33 au total (énumération par solveur, non certifiée) ;
(v) 8 de ces 33 problèmes sont insatisfiables, et \textbf{les 8 sont certifiés} : preuve DRAT produite par un second solveur (CaDiCaL 3.0.1)
et vérifiée par \code{drat-trim}.
\textbf{Aucune classe n'est encore entièrement exclue ; la branche C3 n'est pas tranchée.}
Le document distingue systématiquement preuve, certificat, calcul et résultat négatif, liste les erreurs commises et corrigées en cours de route,
et chiffre ce qu'il faudrait pour aller au bout.
\end{abstract}

\tableofcontents

\section*{Échelle de statut}
\begin{description}[nosep]
\item[\st{D}] \emph{Démontré} : preuve écrite dans ce document, relue ; à revérifier par un humain.
\item[\st{L}] \emph{Littérature} : repris d'une source primaire lue (référence donnée).
\item[\st{C}] \emph{Certifié} : le résultat d'un calcul est accompagné d'un certificat exact (témoin, vecteur de Farkas, preuve DRAT) vérifié par un programme
indépendant du solveur qui l'a produit.
\item[\st{S}] \emph{Calcul par solveur, non certifié} : réponse d'un solveur (CP-SAT, CaDiCaL) sans certificat vérifié.
\item[\st{N}] \emph{Résultat négatif} : une piste testée qui n'élimine rien.
\item[\st{?}] \emph{Ouvert / à faire}.
\end{description}

%==================================================================
\section{Résumé pour le lecteur pressé}
\begin{enumerate}[nosep]
\item \st{L}+\st{D} Un automorphisme d'ordre 3 de $\Gamma$ est sans point fixe et possède \textbf{exactement 6 orbites triangulaires} et 27 orbites indépendantes
(§\ref{sec:tau}). Conséquence annexe : l'exclusion calculée « $f=27$ » du dépôt \code{infinityscroll/conway99-order3-f27} (août 2026) est redondante.
\item \st{C} Le graphe $F$ entre les 6 orbites triangulaires est l'un de 26 graphes sur 156 (§\ref{sec:motifs}).
\item \st{C} Avec les appariements, il reste 125 classes de phases, soit 46 à symétrie près ; on retrouve au passage le « lemme de l'hexagone » (§\ref{sec:phases}).
\item \st{D} Deux lemmes nouveaux au niveau des sommets restreignent les incidences possibles vers les orbites indépendantes (§\ref{sec:l25}).
\item \st{S} Énumération des incidences possibles : complète pour 35 classes sur 46 (1 à 6 « multi-ensembles de supports » par classe, 33 problèmes distincts), incomplète pour 11 (§\ref{sec:enum}).
\item \st{C} Étage quotient (§\ref{sec:stageq}) : sur 33 problèmes, 8 insatisfiables, \textbf{tous certifiés par preuve DRAT vérifiée}
(cf.\ tableau~\ref{sec:tabC}),
20 indécis après 5 minutes, 5 non terminés (arrêtés faute de mémoire). \textbf{Aucune classe n'est encore exclue en entier.}
\item \st{N} Plusieurs pistes n'éliminent rien : règle par colonne, filtres locaux au niveau des classes, condition spectrale sur $C$, solveur sur le modèle complet (§\ref{sec:neg}).
\item Tous les modèles ont été validés sur des graphes réels ($\srg(243,22,1,2)$ de Berlekamp--van Lint--Seidel, graphe de la tour $3\times3$) avant usage ;
six erreurs de programme et deux problèmes de ressources ont été détectés et corrigés avant de produire un résultat faux, et l'audit a corrigé deux
affirmations inexactes de comptes rendus intermédiaires (§\ref{sec:incidents}).
\end{enumerate}

%==================================================================
\section{Cadre, notations, faits de base}\label{sec:cadre}
Soit $\Gamma$ un $\srg(99,14,1,2)$ hypothétique, $A$ sa matrice d'adjacence : $A^2+A=12I+2J$, spectre $\{14^1,3^{54},(-4)^{44}\}$.
Soit $g\in\Aut(\Gamma)$ d'ordre 3.

\paragraph{Absence de point fixe \st{L}.} Makhnev--Minakova (Diskr.\ Mat.\ 16:1, 2004, th.~1) : $\Fix(g)$ est vide ou un triangle.
Behbahani (thèse 2009, th.~4.14) / Behbahani--Lam (2011) : pas de matrice d'orbites pour $p=3$ avec points fixes. Donc $\Fix(g)=\emptyset$ :
33 orbites de taille 3. On note les sommets $(U,a)$, $a\in\Z/3$, avec $g(U,a)=(U,a+1)$.
Une orbite est \emph{triangulaire} si elle induit $K_3$, \emph{indépendante} sinon ($g$ agit transitivement sur les trois paires d'une orbite).
On note $\tau$ le nombre d'orbites triangulaires.

\paragraph{Quotient.} $Q_{UV}=|N(u)\cap V|$ pour $u\in U$. Toutes les orbites ont la même taille, donc $Q$ est symétrique, $Q\mathbf 1=14\cdot\mathbf 1$ et
\[Q^2+Q=12I+6J .\]

\begin{lemma}[entrées du quotient] \st{D}\label{lem:entries}
(a) Si $U$ est triangulaire et $V\neq U$, alors $Q_{UV}\le 1$ ; les arêtes entre $U$ et $V$ forment alors un appariement $g$-équivariant.
(b) Si $U,V$ sont indépendantes, $Q_{UV}\le 2$.
(c) Toute orbite indépendante $U$ a exactement deux orbites $V$ avec $Q_{UV}=2$ (les « doubles » forment un 2-facteur sur les orbites indépendantes).
\end{lemma}
\begin{proof}
(a) Si $Q_{UV}\ge2$, par symétrie un sommet $v\in V$ a deux voisins $u,u'$ dans le triangle $U=\{u,u',u''\}$ ; l'arête $uu'$ aurait deux voisins communs
$u''$ et $v$, contre $\lambda=1$.
(b) Si $u\in U$ était adjacent aux trois sommets de $V$, les sommets non adjacents $u$ et $gu$ auraient trois voisins communs, contre $\mu=2$.
(c) $u$ et $gu$ ne sont pas adjacents, donc ont exactement deux voisins communs. Dans une orbite $W$ avec $Q_{UW}=1$ ($u\sim w$, $gu\sim gw$) : aucun.
Dans une orbite triangulaire : aucun, par (a). Dans $W$ avec $Q_{UW}=2$ ($u\sim w,g^\delta w$, $gu\sim gw,g^{\delta+1}w$, $\delta\in\{1,2\}$) :
exactement un. D'où exactement deux orbites doubles.
\end{proof}

On ordonne les orbites : les $\tau$ triangulaires d'abord. Alors
\[Q=\begin{pmatrix}2I+F & B\\ B^{\mathsf T}& C\end{pmatrix},\qquad F\in\{0,1\}^{\tau\times\tau},\ B\in\{0,1\}^{\tau\times(33-\tau)},\ C\in\{0,1,2\}^{\text{indép.}}.\]

\paragraph{Reformulation exacte (anneau de groupe).} $A$ est une matrice $33\times33$ à coefficients dans $\Z[C_3]=\Z[x]/(x^3-1)$ ; comme
$\Z[C_3]\hookrightarrow\Z\times\Z[\omega]$, l'équation $A^2+A=12I+2J$ équivaut à deux équations : en $x=1$ l'équation du quotient ci-dessus, et en $x=\omega$
une matrice hermitienne $M$ à coefficients d'Eisenstein (entrées $0$, $\omega^k$ pour $Q=1$, $-\omega^k$ pour $Q=2$, diagonale $-1$ ou $0$) avec
$M^2+M=12I$, spectre $\{3^{18},(-4)^{15}\}$ (analogue de la matrice signée $S$ de la branche involution). Cette décomposition guide les « étages » ci-dessous.

%==================================================================
\section{Résultat 1 : exactement six orbites triangulaires}\label{sec:tau}
\begin{proposition} \st{L}+\st{D}
$\tau=6$ : $g$ a six orbites triangulaires et 27 orbites indépendantes.
\end{proposition}
\begin{proof}
Soit $E_3=\frac{4I+A}{7}-\frac{2J}{77}$ le projecteur sur l'espace propre de valeur 3 (entrées : $\tfrac{42}{77}$ diagonale, $\tfrac{9}{77}$ adjacents, $-\tfrac{2}{77}$
non adjacents). Le caractère vaut
$\chi_3(g)=\tr(gE_3)=\sum_v E_3[v,gv]=3\tau\cdot\tfrac{9}{77}-(99-3\tau)\tfrac{2}{77}=\tfrac{3(\tau-6)}{7}$ (les $3\tau$ sommets des orbites triangulaires sont
adjacents à leur image, les autres non). L'intégralité seule donne $\tau\equiv6\pmod 7$, soit $\tau\in\{6,13,20,27\}$.
Ishida (arXiv:2606.29183v2, prop.~4.4 et 8.7) : si $E$ est un idempotent spectral à entrées $p$-entières et $x$ d'ordre $p$, alors
$\tr(xE)=\rg_{\F_p}\bigl(E[\Fix(x)]\bmod p\bigr)$. Ici $p=3$ ; $E_3$ est $3$-entier (dénominateurs $7$ et $77$) ; $\Fix(g)=\emptyset$, donc le rang est nul,
$\chi_3(g)=0$ et $\tau=6$.
Nous avons relu l'argument d'Ishida : l'image de $E$ est un facteur direct du réseau de permutation $\Z_3^{V}$, donc (théorème de Green + Krull--Schmidt sur
$\Z_3$) une somme de réseaux $\Z_3[G/H]$ ; la trace de $x$ compte les facteurs triviaux, que le quotient de Brauer, égal à $E[\Fix(x)]\bmod 3$, compte aussi.
Contrôle de cohérence (calcul à la main) sur un vrai graphe : $\srg(243,22,1,2)$ avec l'involution $x\mapsto-x$ ($p=2$, un point fixe) : trace et rang valent $0$.
\end{proof}
\begin{remark}
Le résultat est dû à Ishida (juillet 2026). Il rend redondante l'exclusion calculée du cas « $f=27$ » par \code{infinityscroll/conway99-order3-f27}
(14 août 2026), dont le README présente encore $f=6,13,20$ comme ouverts. Il lève aussi la réserve du dossier de départ (ChatGPT), qui traitait
« six orbites triangulaires » comme une hypothèse de travail.
\end{remark}

%==================================================================
\section{Résultat 2 : le Gram des six triangles}\label{sec:gram}
\begin{lemma} \st{D}
$G_0:=BB^{\mathsf T}=6I+6J-5F-F^2$. En particulier la somme de la ligne $i$ de $B$ vaut $12-\deg_F(i)$, $G_0\ge0$ entrée par entrée et $G_0\succeq0$.
\end{lemma}
\begin{proof}
Bloc triangulaire de $Q^2+Q=12I+6J$ : $(2I+F)^2+(2I+F)+BB^{\mathsf T}=12I+6J$.
\end{proof}
Si $s\in\{0,1\}^6$ désigne le support d'une colonne de $B$ (les triangles adjacents à une orbite indépendante) et $m_s$ le nombre de colonnes égales à $s$ :
\begin{equation}\label{eq:agg}
\textstyle\sum_s m_s=27,\qquad \sum_s m_s\,ss^{\mathsf T}=G_0,\qquad m_s\in\Z_{\ge0}.
\end{equation}

%==================================================================
\section{Résultat 3 : 26 graphes $F$ possibles sur 156}\label{sec:motifs}
\begin{theorem} \st{C}
Parmi les 156 graphes à 6 sommets (à isomorphisme près), seuls 26 admettent une solution de~\eqref{eq:agg}. Les 130 autres sont exclus :
87 car $G_0$ a une entrée négative, 25 car un mineur principal de $G_0$ est négatif, 18 par un \emph{certificat de Farkas entier}
$y$ : $y\cdot a(s)\ge0$ pour les 64 colonnes $s$ ($a(s)=(1,(s_is_j)_{i\le j})$) et $y\cdot b<0$ ($b=(27,(G_{0,ij})_{i\le j})$), ce qui interdit toute solution
$m\ge0$ \emph{même fractionnaire}.
\end{theorem}
\textbf{Vérification.} Les certificats ont été cherchés avec CP-SAT puis vérifiés en arithmétique entière. Un second vérificateur, écrit indépendamment
(\code{audit\_verify.py}), repart de l'atlas des graphes de \code{networkx}, recalcule $G_0$, les mineurs (rationnels exacts), revérifie les 26 témoins et les
18 vecteurs de Farkas, sans utiliser aucune donnée du dossier ChatGPT : \code{A1: 87/25/44 ; A2: 26 témoins, 18 Farkas}. Le tableau~A donne les 44 graphes à
Gram positif. Les 18 exclusions par Farkas étaient auparavant des « infaisabilités MILP non certifiées » (note ChatGPT, §7).

%==================================================================
\section{Résultat 4 : phases, 125 classes, 46 à symétrie près}\label{sec:phases}
\paragraph{Modèle \st{D}.} Sur une arête $ij$ de $F$, l'appariement s'écrit $(i,a)\sim(j,a+\varepsilon_{ij})$. Soit $P$ la matrice $18\times18$ induite sur les
sommets des triangles. Si $Y$ est la matrice $18\times81$ des adjacences vers les sommets indépendants,
\[YY^{\mathsf T}=R:=12I+2J-P^2-P .\]
Une orbite indépendante de support $s$ et de phases $p$ (sommet $(Z,b)\sim(i,b+p_i)$) contribue exactement $1$ à $R[(i,0),(j,p_j-p_i)]$.
D'où un problème linéaire en entiers exact (« niveau 2 ») sur les 1366 colonnes phasées. Changer l'étiquetage d'un triangle ($g_i\in\Z/3$) change
$\varepsilon_{ij}$ en $\varepsilon_{ij}+g_j-g_i$ : les classes utiles sont celles de l'espace des cycles, $3^{\beta_1(F)}$ par graphe.

\begin{theorem} \st{C}
Pour les 26 graphes, il y a 318 classes de phases. Pour 191, $R$ a une entrée négative (certificat : l'entrée). Pour 2, la décomposition est impossible :
ce sont exactement le $C_6$ d'holonomie $1$ et $2$ (lemme de l'hexagone, démontré dans le dossier de départ et revérifié). Les 125 autres admettent une
décomposition exacte (témoin vérifié). Modulo $\Aut(F)\times$ jauge $\times$ conjugaison ($g\mapsto g^{-1}$, soit $\varepsilon\mapsto-\varepsilon$), il reste
\textbf{46 classes} (tableau~B).
\end{theorem}
\textbf{Vérification.} Recalcul indépendant (\code{audit\_verify.py}, partie B) : 318 / 191 / 125 témoins exacts / 2 cas $C_6$ d'holonomie $2$ et $1$ ;
partie C : 125 classes $\to$ 46 orbites, avec une autre forme canonique. Le fait que notre modèle retrouve exactement le lemme de l'hexagone est une
validation croisée du modèle par un résultat démontré.

%==================================================================
\section{Résultat 5 : deux lemmes au niveau des sommets}\label{sec:l25}
\begin{lemma}[filtre par type] \st{D}\label{lem:type}
Soit $Z$ une orbite indépendante de support $s$ et phases $p$. Pour tout triangle $i$ et tout $b\in\Z/3$, en notant $e_{ij}$ la phase de l'arête $ij$ orientée
de $i$ vers $j$ :
$\#\{j\in N_F(i)\cap s : b+p_j\equiv e_{ij}\}\le 1$ si $i\in s$, et $\le2$ si $i\notin s$.
\end{lemma}
\begin{proof}
Voisins communs de $u=(i,0)$ et $v=(Z,b)$ parmi les sommets des triangles. Dans $T_j$ ($j\in N_F(i)$) : le voisin de $u$ est $(j,e_{ij})$, celui de $v$ est
$(j,b+p_j)$ si $j\in s$. Dans $T_i$ : si $i\in s$ et $b+p_i\equiv0$, $u\sim v$ et $T_i$ n'apporte rien ; si $i\in s$ et $b+p_i\not\equiv0$, $u\not\sim v$
et $T_i$ apporte un voisin commun ; si $i\notin s$, $u\not\sim v$ et rien. On conclut avec $\lambda=1$, $\mu=2$.
\end{proof}
\begin{lemma}[triplets] \st{D}\label{lem:triple}
Pour tout triplet de triangles $\{i,j,l\}$ et tout couple de phases relatives $r\in(\Z/3)^2$, au plus \emph{une} orbite indépendante (comptée avec
multiplicité) contient $\{i,j,l\}$ avec les phases relatives $(p_j-p_i,p_l-p_i)=r$.
\end{lemma}
\begin{proof}
Si deux orbites $Y\neq Z$ le faisaient, $p^Y-p^Z$ serait constant ($=c$) sur le triplet ; les sommets $(Y,0)$ et $(Z,c)$ auraient alors au moins trois voisins
communs dans les triangles, contre $\lambda,\mu\le2$.
\end{proof}
\textbf{Effet.} Ces lemmes n'éliminent aucune des 125 classes \st{N}, mais ils réduisent fortement l'ensemble des incidences $B$ possibles, et
\textbf{les exclusions de l'étage quotient en dépendent}. Contrôle d'audit : sans ces filtres, l'énumération (niveau 2 seul) donne des sur-ensembles
(classes 0, 2, 3, 11, 18, 19, 26, 34, 35 : 6/6, 9/5, 5/3, 6/2, 4/2, 8/2, 10/3, 10/3, 6/3 supports), sans jamais perdre un support filtré.
Ces lemmes n'ont pas pu être testés sur un graphe réel (aucun graphe connu n'a la structure « 6 triangles ») : \textbf{leur relecture humaine est prioritaire}.
Alternative robuste : passer aussi les supports supplémentaires dans l'étage quotient.

%==================================================================
\section{Résultat 6 : énumération des incidences $B$}\label{sec:enum}
\st{S} Pour chacune des 46 classes, CP-SAT énumère toutes les solutions du niveau 2.5 (niveau 2 + lemmes~\ref{lem:type} et~\ref{lem:triple}), puis on garde
les \emph{multi-ensembles de supports} (seule donnée utile à l'étage quotient). Limite de 120\,s par classe. Résultat (tableau~B) :
\textbf{35 classes complètement énumérées}, 1 à 6 multi-ensembles de supports chacune, soit \textbf{33 problèmes distincts} (les classes d'un même graphe
partagent leurs supports) ; \textbf{11 classes incomplètes} (graphes $F$ peu denses, masques 0, 28, 8290, 9313, 13452, 16384, 16388, 16420, 26896, 29448,
30992 : 1\,163 supports trouvés avant la limite, jusqu'à 560 pour une seule classe).
\emph{Correction d'audit} : des comptes rendus intermédiaires annonçaient « 33 classes complètes, 13 incomplètes » ; 33 est le nombre de problèmes distincts,
les classes complètes sont 35. Aucun calcul n'en dépendait (le balayage sélectionne les classes par leur drapeau « complet »).
Statut : la complétude est une affirmation de CP-SAT, non certifiée ; elle a été recoupée, pour 9 classes, par une énumération indépendante sans les filtres.

%==================================================================
\section{Résultat 7 : l'étage quotient}\label{sec:stageq}
\paragraph{Problème.} Pour un graphe $F$ et un multi-ensemble de supports (colonnes de $B$ affectées aux orbites 6--32), existe-t-il
$C\in\{0,1,2\}^{27\times27}$ symétrique, diagonale nulle, avec exactement deux entrées $2$ par ligne et $Q^2+Q=12I+6J$, $Q\mathbf 1=14\cdot\mathbf1$ ?
C'est une relaxation exacte (conséquence nécessaire) ; si c'est impossible pour \emph{tous} les supports d'une classe, la classe est exclue.

\paragraph{Encodage et validation.} Entrées $C_{UV}=u_{UV}+w_{UV}$ ($w\Rightarrow u$), produits par portes ET (Tseitin), sommes par cardinalités (totalizer de
PySAT), littéraux répétés remplacés par des copies équivalentes. Validation \emph{avant} usage : (1) le vrai quotient de $\srg(243,22,1,2)$ sous une
translation d'ordre 3 (81 orbites triangulaires, puis 81 indépendantes) satisfait toutes les contraintes générées, et une modification est détectée
(43 violations) ; (2) le chemin CNF complet sur le graphe de la tour $3\times3$ : SAT avec vérification numérique pour $k=4$, UNSAT pour $k=5$.

\begin{lemma}[cassage de symétrie des colonnes jumelles] \st{D}\label{lem:lex}
Deux orbites indépendantes de même support sont interchangeables. Si une solution existe, il en existe une où, pour tous jumeaux $i<j$,
la ligne $i$ de $Q$ est lexicographiquement $\le$ la ligne $j$, en comparant les colonnes $W\notin\{i,j\}$ dans l'ordre croissant.
\end{lemma}
\begin{proof}
Soit $A$ la solution minimale pour l'ordre lexicographique « ligne par ligne » dans son orbite sous le groupe $G$ des permutations de jumeaux
(orbite finie, formée de solutions). Supposons $A[i]\succ A[j]$ (colonnes hors $\{i,j\}$) et soit $k$ la première colonne où elles diffèrent :
$A_{ik}>A_{jk}$ ; $k$ est une orbite indépendante (les colonnes des triangles sont égales pour des jumeaux). Soit $B=\tau A\tau$, $\tau=(i\,j)\in G$.
Si $k<i$ : les lignes $p<k$ sont inchangées (car $A_{pi}=A_{ip}=A_{jp}=A_{pj}$ pour $p<k$) et, ligne $k$, la première colonne modifiée est $i$ où
$B_{ki}=A_{jk}<A_{ik}=A_{ki}$ : $B\prec A$. Si $k>i$ : les lignes $p<i$ sont inchangées, et la ligne $i$ de $B$ coïncide avec celle de $A$ avant la colonne $k$
(diagonales nulles, $B_{ij}=A_{ij}$) puis $B_{ik}=A_{jk}<A_{ik}$ : $B\prec A$. Contradiction dans les deux cas.
Nous n'imposons la contrainte qu'entre jumeaux consécutifs, ce qui est plus faible.
\end{proof}
C'est l'adaptation à notre cadre (sous-groupe, entrées $\{0,1,2\}$, parties fixées) du cassage de Codish, Miller, Prosser et Stuckey (IJCAI 2013).

\paragraph{Résultats.} \st{S}/\st{C} Tableau~C : 8 problèmes UNSAT d'après CaDiCaL (pysat, 16 à 100\,s), 20 indécis après 5\,min, 5 non terminés
(processus tués faute de mémoire sur le PC portable). Pour obtenir des \emph{certificats}, chaque UNSAT est rejoué avec CaDiCaL 3.0.1 natif (second solveur),
qui écrit une preuve DRAT, vérifiée par \code{drat-trim} ; la CNF exacte est identifiée par son empreinte SHA-256.
\textbf{Les 8 UNSAT sont certifiés} (\code{certs/SUMMARY.txt}) : CaDiCaL 3.0.1 confirme chaque fois UNSAT (37 à 353\,s) et \code{drat-trim} répond
« s VERIFIED » (36 à 210\,s), sur 1 cœur et moins de 3\,Go ; par exemple m16927\_s0023 : 1\,424 lemmes sur 1,1 million dans le cœur de la preuve.
Chaque CNF (environ 607\,000 variables, 4,7 millions de clauses) est identifiée par son SHA-256 (\code{certs/*.gen.json}) et regénérable à l'identique.
\emph{Portée exacte du certificat} : il établit que la CNF est insatisfiable, indépendamment du solveur. Qu'elle encode fidèlement le problème mathématique
repose sur la validation du générateur (graphes réels) et sur le lemme~\ref{lem:lex} ; que ces problèmes épuisent une classe repose en plus sur les
lemmes~\ref{lem:type}--\ref{lem:triple} et sur la complétude de l'énumération (§\ref{sec:enum}).

\paragraph{Ce que cela implique aujourd'hui.} Aucune classe n'est exclue en entier. Les plus proches : graphe 29481 (classes 26, 34, 35) : 2 supports sur 3
UNSAT ; graphe 5033 (classe 2) : 3 sur 5 ; graphe 16927 (classe 0) : 2 sur 6.

%==================================================================
\section{Résultats négatifs et pistes infructueuses}\label{sec:neg}
\begin{itemize}[nosep]
\item \st{N} Règle par colonne $3s_i+(Fs)_i\le6$ (bloc triangle/indépendant du quotient) : n'élimine aucun graphe.
\item \st{N} Lemmes~\ref{lem:type} et~\ref{lem:triple} : n'éliminent aucune des 125 classes (mais filtrent $B$, cf.\ §\ref{sec:l25}).
\item \st{N} \textbf{Condition spectrale} \st{D} : de $BC=6J-(3I+F)B$ et $C\mathbf1=14\mathbf1-B^{\mathsf T}\mathbf1$ on tire $CM=MK$ avec $M=[B^{\mathsf T}|\mathbf1]$ et
$K=\begin{pmatrix}-(3I+F)&-\mathbf1\\6\cdot\mathbf1^{\mathsf T}&14\end{pmatrix}$ ; $C$ préserve $W=\operatorname{Im}M$ et $W^\perp$, où $C^2+C=12I$.
Le spectre de $C$ est donc entièrement déterminé par $(F,\text{supports})$ ; conditions nécessaires : valeurs propres réelles, multiplicités entières,
$\tr C^3\ge0$ et $\equiv0\pmod6$. La dérivation est validée ($\tr C^2$ prédit = $\tr C^2$ combinatoire dans les 33 cas), mais \textbf{les 33 cas passent} :
la condition n'élimine rien. (Un faux « échec » est apparu d'abord : cf.\ §\ref{sec:incidents}.)
\item \st{N} Modèle complet au niveau des sommets (triangles + 27 orbites indépendantes + phases), CP-SAT puis CaDiCaL, avec ou sans $B$ fixé :
indécis en 2 à 10 minutes sur la classe la plus contrainte.
\end{itemize}

%==================================================================
\section{Validation, erreurs détectées et corrigées}\label{sec:incidents}
Aucun des incidents ci-dessous n'a affecté un résultat présenté ; chacun a été détecté par une vérification avant d'être rapporté.
\begin{longtable}{p{4.3cm}p{5.6cm}p{5.6cm}}
\toprule Incident & Détection & Effet / correction\\\midrule
Portes ET mal encodées dans le constructeur CP-SAT (\code{OnlyEnforceIf(g)} au lieu de \code{g.Not()}) & « UNSAT » en 2\,s sans conflit, noyau d'une seule
contrainte manifestement satisfiable & faux UNSAT \emph{rejeté} ; constructeur corrigé puis validé sur la tour $3\times3$ (SAT/UNSAT attendus)\\
Interruption de CaDiCaL via PySAT inopérante & processus toujours actif après 10\,min pour une limite de 2\,min & remplacée par un plafond au niveau du processus\\
Clause vide non gérée par PySAT/CaDiCaL & plantage du processus & rapporté comme « indécis », jamais comme UNSAT ; remplacée par une contradiction explicite\\
Littéraux répétés dans les cardinalités & revue du code (risque, non observé) & remplacés par des copies équivalentes ; modèles revalidés\\
Comptage de Sturm (racines \emph{distinctes}) & « valeurs propres non réelles » sur 7 graphes entiers : trop beau & faux positif \emph{rejeté} ; test sur la partie sans
facteur carré\\
Retour chariot dans la sortie de \code{drat-trim} & « NOT VERIFIED » alors que le journal dit « s VERIFIED » & erreur dans le sens prudent ; lecture corrigée et testée\\
Recherche exhaustive Python trop lente, 6 processus à 100\,\% & gel ressenti du PC & arrêtée ; remplacée par CP-SAT dans Docker plafonné\\
Mémoire insuffisante (8 puis 5 CaDiCaL en parallèle) & processus tués & 5 problèmes non terminés ; à relancer sur une machine à 32\,Go\\
Comptes rendus intermédiaires : « 33 classes complètes, 13 incomplètes » & recomptage pendant l'audit & erreur de texte seulement : 35 et 11 (33 = nombre de
problèmes) ; corrigé\\
Affirmation intermédiaire « le niveau 2.5 n'élimine rien » & comparaison des énumérations avec et sans filtres & vrai au niveau des classes, faux au niveau des $B$ ;
les exclusions dépendent donc des lemmes~\ref{lem:type}--\ref{lem:triple} (précisé §\ref{sec:l25})\\
\bottomrule
\end{longtable}

%==================================================================
\section{Ce que cela signifie, et ce que cela ne signifie pas}
\begin{itemize}[nosep]
\item \textbf{Rien n'est tranché.} La branche C3 reste ouverte ; aucune classe n'est exclue en entier.
\item \textbf{Ce qui est acquis} (\st{D}/\st{C}) : une réduction exacte et vérifiable de la branche C3 à 46 classes de phases, dont 35 se ramènent chacune à 1--6
problèmes quotient explicites ; 8 de ces problèmes sont insatisfiables, avec preuves DRAT vérifiées.
\item \textbf{Si} tous les problèmes restants (et, pour les survivants, l'étage des phases sur $C$) étaient insatisfiables, $\Gamma$ n'aurait pas
d'automorphisme d'ordre 3, donc $|\Aut(\Gamma)|\in\{1,2\}$. Joint à la branche involution (dépôt \code{conway99-involution}), cela donnerait
$|\Aut(\Gamma)|=1$. Cela ne tranche pas l'existence de $\Gamma$.
\item \textbf{Nouveauté} : non vérifiée. Le dépôt public \code{ibenarb/conway99-research} (« types O3 ») et \code{Lundle22/conway-99-cyclic-lifts} contiennent
peut-être des résultats équivalents (classification des graphes $F$, réduction par phases). La priorité n'est revendiquée pour rien.
\end{itemize}

%==================================================================
\section{Ce qu'il faut pour aller au bout}\label{sec:suite}
\subsection*{Vérifications humaines (par ordre d'importance)}
\begin{enumerate}[nosep]
\item Lemme~\ref{lem:entries} et la forme de $Q$ (15 min).
\item Lemmes~\ref{lem:type} et~\ref{lem:triple} : \emph{les exclusions en dépendent} (20 min).
\item Lemme~\ref{lem:lex} (cassage de symétrie) : les UNSAT en dépendent (20 min).
\item Argument $\tau=6$ (Ishida) (15 min) et le modèle de phases (§\ref{sec:phases}) (20 min).
\item Rejouer \code{check\_all.py} et \code{audit\_verify.py} (bibliothèque standard + networkx ; quelques minutes).
\end{enumerate}
\subsection*{Calculs}
\begin{enumerate}[nosep]
\item \emph{Fait} : certification DRAT des 8 UNSAT ($\le 10$\,min par cas, 1 cœur, $<3$\,Go). Tout nouvel UNSAT se certifie de la même façon.
\item Relancer les 5 problèmes non terminés, puis les 20 indécis avec un budget croissant (1\,h, puis 24\,h par problème). Coût inconnu : c'est la mesure
à faire. Un solveur CaDiCaL occupe environ 1\,Go ; une machine à 32\,Go et 8--16 cœurs en fait tourner 12--16 à la fois.
\item Robustesse : passer aussi à l'étage quotient les supports « sans filtres » des classes proches de l'exclusion (au plus 35 problèmes de plus).
\item Compléter l'énumération des 11 classes à $F$ peu dense (formulation directe au niveau des supports ; plusieurs milliers de problèmes possibles).
\item Pour tout problème quotient qui serait satisfiable : étage des phases sur $C$ (matrice d'Eisenstein $M^2+M=12I$).
\item Rendre l'énumération des $B$ certifiable (par exemple énumération par clauses bloquantes avec preuve DRAT de la dernière insatisfiabilité).
\end{enumerate}
\subsection*{Matériel}
La machine de calcul du projet (Ryzen 7 7800X3D, 32\,Go) suffit pour les points 1--3 ; les points 4--5 relèvent d'un serveur de calcul (dizaines de cœurs,
jours de calcul) ou d'un accès à un cluster universitaire.

%==================================================================
\section{Complément de la nuit du 22 au 23 septembre 2026}\label{sec:nuit}
\emph{Travail autonome (l'opérateur dormait). Même échelle de statut. Détail horodaté : \code{NUIT\_JOURNAL.md}.}

\subsection{Recherche bibliographique ciblée}
\begin{itemize}[nosep]
\item \st{L} \code{ibenarb/conway99-research} (dépôt actif, lu le 23/09) : même quotient $Q=2D_T+S+2L$ ; démontre $\tau\equiv0\pmod3$ par un comptage
élégant (les $231$ triangles du graphe ; les triangles fixés par $g$ sont exactement les $\tau$ orbites triangulaires) ; exclut les 4-cycles du 2-facteur $L$
et démontre le « lemme B » (une composante triangulaire $\{a,b,c\}$ de $L$ reçoit exactement une arête simple de chaque autre orbite).
Découpage par type de cycles de $L$ : 103 types « ouverts » de leur point de vue ; avec Ishida ($\tau=6$), leurs 2 types $\tau=27$ sont vides.
Coût mesuré de leur méthode : l'exclusion d'\emph{un} type, $(6,3^7)$, a demandé 656 sous-cas et 652\,Go de preuves LRAT.
\item \st{L} \code{Lundle22/conway-99-cyclic-lifts} : sous l'équation du quotient, l'équation complète équivaut à la congruence $W^2+W-12I\equiv0\pmod{\pi^3}$
($\pi=1-\omega$) sur la matrice d'Eisenstein, soit une équation linéaire et une quadratique sur $\F_3$ (lemme clé revérifié : trois entiers $\ge0$ de somme 6
avec $\pi^3\mid c_0+c_1\omega+c_2\omega^2$ valent tous 2). Utile pour l'étage des phases.
\item \st{L} Thakkar--Severini (arXiv:2608.11211) : modèle d'orbites générique ; ordre 3 sans point fixe indécis en 30\,min. SAT Modulo Symmetries :
non adapté (pas de multigraphe ni de partition imposée). « Smart cubing » : gains $\times2$--$3$.
\end{itemize}

\subsection{La relaxation linéaire LQ}
\begin{lemma}[bloc triangles $\times$ indépendantes] \st{D}\label{lem:LQ}
Pour toute orbite indépendante $Z$ de support $s_Z$ : $\sum_{Y\neq Z}C_{YZ}\,s_Y=t_Z:=6\cdot\mathbf1-3s_Z-Fs_Z$, $\sum_Y C_{YZ}=14-|s_Z|$, et exactement deux
$Y$ ont $C_{YZ}=2$ ; en particulier $s_Y(i)\,C_{YZ}\le t_Z(i)$.
\end{lemma}
\begin{proof}
Bloc $(i,Z)$ de $Q^2+Q=12I+6J$ : $3s_Z(i)+(Fs_Z)_i+\sum_Y s_Y(i)C_{YZ}=6$ ; somme de ligne de $Q$ ; lemme~\ref{lem:entries}(c).
\end{proof}
\paragraph{Résultats.} Sur les 33 problèmes quotient : \textbf{15 infaisables} pour LQ, dont les 8 déjà certifiés et \textbf{7 nouveaux}
(m04776\_s0000, s0002, s0003 ; m04905\_s0004, s0005 ; m05033\_s0008 ; m15426\_s0017). \textbf{13 des 15 le sont par une seule ligne} : une orbite $Z$ n'admet
aucun voisinage compatible ; certificat = recherche exhaustive sur une ligne (Python pur, \code{lq\_local.py}), recoupée par un modèle CP-SAT indépendant
sur les 891 lignes des 33 problèmes (0 désaccord). Les 2 autres (m16927\_s0018, s0023) demandent la symétrie globale et sont certifiés DRAT.
La variante la plus faible (sans lemme B, sans règle des 4-cycles, sans borne sur les doubles) suffit dans les 15 cas.

\paragraph{Contenu linéaire épuisé.} \st{D} Toute somme de $C^2$ contre un vecteur de $W=\operatorname{Im}[B^{\mathsf T}|\mathbf1]$ se réduit, via le
lemme~\ref{lem:LQ}, à une identité : LQ contient tout le contenu linéaire. Le reste est quadratique, dans $W^\perp$ ($C^2+C=12I$ y restreint).

\subsection{Classes clairsemées : énumération directe des supports}
Au lieu d'énumérer les $B$ avec leurs phases (explosif pour $F$ peu dense), on énumère les multi-ensembles de supports compatibles avec LQ, graphe par graphe
(\code{enum\_supports\_lq.py} : supports triés + $C$ + contraintes linéaires ; chaque solution est interdite puis on recommence ; fin sur « infaisable »
$\Rightarrow$ liste complète).
\begin{itemize}[nosep]
\item $F=K_{1,3}$ (masque 28, classe 42, auparavant 637 $B$ phasés sans fin d'énumération) : \textbf{1 seul multi-ensemble}, liste complète en 59\,s.
\item $F=3K_2$ (masque 16420, classe 27, auparavant 1\,817 $B$ phasés) : \textbf{4 multi-ensembles}, liste complète en 102\,s ; cohérent avec l'ancienne
liste partielle (les 2 autres supports qu'elle contenait échouent à LQ).
\item Masque 8290 (classe 25) : 2 multi-ensembles ; masque 9313 (classe 32) : 2 ; masque 13452 (classe 5) : 1 ; masque 29448 (classe 6) : 1.
Chaque fois, les multi-ensembles de l'ancienne liste partielle absents de la nouvelle échouent bien à LQ (\code{compare\_supports.py}).
\end{itemize}
Statut \st{S} : complétude affirmée par CP-SAT.

\paragraph{Test de phases sur un multi-ensemble de supports.} \code{support\_phase\_check.py} : existe-t-il un $B$ phasé de niveau 2.5 (résidu exact,
filtre de type, condition des triplets) d'une classe de $F$ dont les supports sont exactement ce multi-ensemble ? Si aucune classe ne le réalise,
le multi-ensemble est exclu. C'est valide avec les représentants de classes : tout graphe se réétiquette pour que sa classe soit un représentant, et son
multi-ensemble dans cet étiquetage est alors réalisé. \textbf{m08290\_sL1 et m09313\_sL1 sont exclus} ; certificat : vecteur de Farkas entier
(\code{farkas\_phase.py}), revérifié en arithmétique exacte \st{C}. Bilan des classes clairsemées : classe 42 $\to$ 1 problème quotient, 27 $\to$ 4,
25 $\to$ 1, 32 $\to$ 1, 5 $\to$ 1, 6 $\to$ 1.

\subsection{Accélération de l'étage Q : encodage compact et symétrie complète}
\paragraph{Encodage compact \st{D}+\st{C}.} L'ancien générateur encodait chaque contrainte « $\sum\ell=r$ » par \code{CardEnc.equals} (totalizer) ;
la moitié « $\ge r$ » y devient « au plus $n-r$ » sur les négations, soit un totalizer presque complet en $O(n^2)$ pour les $\approx100$ littéraux
d'une équation de Gram où $r\le6$. \code{lean\_encoding.py} utilise \emph{un} totalizer bidirectionnel dont les sorties sont tronquées à $r+1$
($o_j\Leftrightarrow\sum\ge j$), tout aussi arc-cohérent, en $O(nr^2)$ clauses ; variables et contraintes sont inchangées.
Validation (\code{validate\_lean.py}, PASS) : 5\,116 affectations exhaustives ($n\le8$, tous $r$) ; quotients réels BvLS (81 T et 81 I, lignes
libres : SAT et évalué ; entrée corrompue : UNSAT, pour les deux encodages) ; graphe tour $3\times3$ ; deux instances certifiées UNSAT.
\begin{center}\begin{tabular}{lrrrr}
instance & clauses (ancien) & clauses (compact) & temps ancien (natif) & temps compact\\\hline
m05033\_s0009 & 4\,727\,236 & 1\,117\,843 & 60\,s & 2{,}8\,s\\
m29481\_s0030 & 4\,721\,852 & 1\,122\,170 & 260\,s & 8{,}0\,s
\end{tabular}\end{center}

\paragraph{Symétrie complète \st{D}.} Jusqu'ici seules les orbites « jumelles » (supports identiques) étaient échangées. Une instance
$(F,\{s_Y\})$ admet aussi les permutations $\pi$ des 6 triangles avec $\pi(F)=F$ et $\pi(\{s_Y\})=\{s_Y\}$, prolongées aux orbites indépendantes :
71 non triviales pour m23221\_s0027, 23 pour m16420\_sL*, 11 pour m21041, m23431, m00028\_sL0. On impose $\mathrm{ch}(Q)\le_{\mathrm{lex}}\mathrm{ch}(Q^\sigma)$
pour chacune, où $\mathrm{ch}$ lit la matrice $33\times33$ complète ligne par ligne (ordre $0<1<2$).
\begin{lemma}\label{lem:symlex}
Soit $Q_0$ le minimum lexicographique de l'orbite d'une solution sous le groupe engendré par ces $\sigma$ et les transpositions de jumelles. Alors
$Q_0$ satisfait les contraintes $\sigma$ et les contraintes jumelles (ligne$_i\le_{\mathrm{lex}}$ligne$_j$ hors colonnes $i,j$, pour $i<j$ jumelles).
\end{lemma}
\begin{proof}
Contraintes $\sigma$ : minimalité. Jumelles : sinon, soit $c$ la première colonne ($\notin\{i,j\}$) où les lignes $i,j$ diffèrent, avec $Q_0[i][c]>Q_0[j][c]$.
Pour $k<c$, $k\notin\{i,j\}$ : $Q_0[k][i]=Q_0[i][k]=Q_0[j][k]=Q_0[k][j]$. Si $c<i$, la transposition $(ij)$ laisse les lignes $<c$ inchangées
et rend la ligne $c$ plus petite en colonne $i$. Si $c>i$, elle laisse les lignes $<i$ inchangées et la ligne $i$ diffère d'abord en $c$, plus petite
(les diagonales coïncident : même type). Contradiction avec la minimalité.
\end{proof}
Validation (\code{validate\_symlex.py}, PASS, avec les fonctions de production) : 60 petits modèles énumérés exhaustivement, groupe complet des données
fixées calculé par force brute ; 4\,044 orbites de solutions, \textbf{aucune perdue}, et le minimum lex de chaque orbite toujours présent ;
36 modèles ont des symétries $\sigma$ actives et des solutions (jusqu'à 13\,950 solutions, 3\,558 orbites). Témoin : la version relâchée de m23221\_s0027
(équations de Gram indépendantes retirées) reste SAT avec les 71 contraintes.

\paragraph{Mémoire.} Solveur alimenté clause par clause depuis le disque (\code{solve\_stream.py}) : 200--350\,Mo par instance au lieu de
$\approx1{,}4$\,Go, d'où 6 solveurs en parallèle au lieu de 2.

\subsection{Premier résultat de l'étage Q accéléré}
\textbf{m16420\_sL0 est UNSAT} \st{C} : CNF compacte avec symétrie complète (sha256 \code{a706468daf92ee33}), UNSAT en 2\,273\,s (CaDiCaL 1.9.5) puis
3\,515\,s (CaDiCaL 3.0.1 natif, autre version) ; preuve DRAT binaire de 129\,Mo vérifiée par drat-trim (2\,852\,s ; 587\,320 lemmes dans le noyau,
aucun lemme RAT). m16420\_sL3 lui est isomorphe. La validité repose en plus sur le lemme~\ref{lem:symlex} et sur la validation de l'encodage.

\subsection{Cube-and-conquer}
Les autres instances ne concluent pas en une heure (étage Q monolithique). Découpage : chaque orbite indépendante a exactement deux doubles
(lemme~\ref{lem:entries}(c)), donc pour une orbite $Z$ fixée, les 325 « cubes » $\{w_{ZA}=w_{ZB}=1\}$ couvrent toutes les solutions. Un cube
récalcitrant se redécoupe de la même façon : une orbite $V$ n'ayant qu'un double supposé (partenaire $P$) a son second double $W\notin\{V,P\}$,
d'où 25 sous-cubes qui couvrent le nœud (profondeur $\le2$). Résolution \emph{incrémentale} (\code{cube\_q\_inc.py}) : un seul solveur par
instance, chaque cube résolu sous hypothèses avec un budget de conflits (5\,000, puis 50\,000, puis découpage) ; les clauses apprises, conséquences
de la CNF de base, servent d'un cube à l'autre.
\paragraph{Validations.} \code{validate\_inc.py} (PASS) : sur le témoin relâché (m23221\_s0027 sans les équations de Gram indépendantes), 325 cubes,
verdicts identiques entre l'incrémental et un solveur neuf avec clauses unitaires (29 SAT, 296 UNSAT, modèles revérifiés clause par clause).
Le gain vient surtout de ce qu'on ne recharge plus 1,4\,M de clauses par cube : la plupart des cubes se résolvent en quelques millisecondes.
\code{validate\_split.py} (PASS) : sur les 60 derniers cubes du témoin avec un budget d'un conflit, 11 cubes tranchés par découpage (7 SAT trouvés
au fond de l'arbre, 1 UNSAT), 0 désaccord, 2\,819 nœuds.

\subsection{Doublons par isomorphisme}
Deux problèmes quotient $(F,\mu)$ et $(F,\pi(\mu))$ avec $\pi\in\operatorname{Aut}(F)$ ont la même réponse. Test systématique sur les 44 fichiers
d'instances : m04776\_s0000$\sim$s0002, m04905\_s0004$\sim$s0005, m16927\_s0018$\sim$s0023, m29481\_s0030$\sim$s0031 (morts des deux côtés),
m16927\_s0019$\sim$s0022, m16927\_s0020$\sim$s0021 (ouverts), m16420\_sL0$\sim$sL3, sL1$\sim$sL2. Les 33 instances d'origine ne font que
\textbf{27 classes d'isomorphisme}, et les 18 ouvertes \textbf{16}. Contrôle de cohérence : chaque paire a le même statut (LQ comme étage Q).

\subsection{Énumération des supports avec les phases}
\code{enum\_supports\_phased.py} : pour une classe $c$, un seul modèle CP-SAT contient le modèle de supports (LQ) et un $B$ phasé de niveau 2.5
de la classe, liés par « nombre de colonnes phasées de support $v$ = nombre d'orbites de support $v$ ». On bloque chaque multi-ensemble trouvé ;
« infaisable » à la fin donne la liste complète des multi-ensembles compatibles LQ \emph{et} réalisables par la classe. Validation
(\code{validate\_enum\_phased.py}, PASS) : sur les classes 0, 15, 30 (listes connues par l'énumération phasée complète) et 25, 32, 42
(clairsemées), ensembles identiques, listes complètes en 17--296\,s.
Variante négative \st{N} : un lex conditionnel des lignes de $C$ pour les supports égaux ralentit CP-SAT (masque 28 : rien en 600\,s au lieu de 135\,s).

\subsection{Pistes testées sans effet \st{N}}
\begin{itemize}[nosep]
\item LQ au niveau d'un graphe $F$ entier (supports libres) : aucun $F$ exclu (21 faisables, 5 indécis en 5\,min).
\item Version « phases » locale du lemme~\ref{lem:LQ} (cible $6\times3$ par orbite ; conventions validées sur 14\,580 comparaisons) : aucune élimination sur 15 classes.
\item Version « phases » globale (\code{phase\_global.py}, validée par contrôles plantés, jamais de faux négatif) : indécise en 60\,s sur les classes
symétriques 36--41.
\item Générateur ligne par ligne : 35\,000 à 2 millions de lignes possibles par orbite ; non viable sans élagage quadratique fort.
\item Balayage de la version globale sur 145 $B$ phasés à 15\,s chacun : 0 faisable, 0 exclu, tout indécis.
\item Relaxation LQB (LQ + inégalité de budget) : indécise en CP-SAT sur toutes les instances testées.
\item Condition spectrale par factorisation de $\chi_M(x^2+x)$ : déjà contenue dans le test spectral de l'audit ($CM=MK$ fixe tout le spectre de $C$).
\item Solutions invariantes (ansatz) pour m23221\_s0027 : aucune sous 8 types de symétries (6 UNSAT, 2 indécis à 900\,s). Rien n'indique que l'étage Q
y soit satisfiable.
\item Les lemmes « pas de $C_4$ dans $L$ » et « lemme B » sont \emph{impliqués} par l'étage Q (Gram $\ge8>6$, resp.\ Gram saturé puis sommes de lignes).
\end{itemize}

\subsection{Bilan au matin du 23 septembre (07\,h\,10) et suite}
\paragraph{Acquis de la nuit.}
\begin{itemize}[nosep]
\item \st{C} m16420\_sL0 (et son isomorphe sL3) : UNSAT, preuve DRAT vérifiée.
\item \st{S} m00028\_sL0 : UNSAT par cube-and-conquer (325/325 cubes, dont 27 fermés par découpage, 675 sous-cubes) ; c'était l'unique instance de la
classe 42 ($F=K_{1,3}$), qui est donc exclue à l'étage quotient. Certification DRAT cube par cube à faire.
\item Instances d'origine : 27 classes d'isomorphisme sur 33 ; ouvertes : 16 sur 18.
\item \st{S} Classe 4 (masque 30992) : liste complète avec phases, 96 multi-ensembles, 18 à isomorphisme près (m30992\_sP0--sP17), pas encore passées à l'étage Q.
\end{itemize}
\paragraph{Instances en cours (cubes UNSAT sur 325, aucun cube SAT).}
\begin{center}\begin{tabular}{lr|lr|lr}
m16420\_sL1 & 316 & m21041\_s0025 & 313 & m23221\_s0027 & 313\\
m23431\_s0028 & 309 & m06717\_s0014 & 307 & m20723\_s0024 & 306\\
m04907\_s0007 & 301 & m12295\_s0015 & 297 & m05033\_s0012 & 291\\
m29481\_s0032 & 288 & & & &
\end{tabular}\end{center}
Pas encore commencées : m04776\_s0001, m04905\_s0006, m15426\_s0016, m16927\_s0019, s0020, m21045\_s0026, m26850\_s0029, m08290\_sL0, m09313\_sL0,
m13452\_sL0, m29448\_sL0, m30992\_sP0--sP17.
\paragraph{Classes encore incomplètes.} 29 ($F=K_{1,5}$ : au moins 16 multi-ensembles avec phases, énumération en cours), 43 (2 trouvés, puis
indécis à 3\,600\,s), 44 et 45 (pas encore traitées avec les phases ; $F$ à une arête et $F$ vide, énumération sans phases indécise).
\paragraph{Coût restant.} Les passes à 5\,000 puis 50\,000 conflits règlent 85 à 95\,\% des cubes en quelques minutes ; chaque cube dur restant
coûte environ 30--45\,min de CPU par découpage. Il en reste de l'ordre de 300 à 400 sur les instances en cours, et les instances non commencées
s'y ajoutent : plusieurs centaines d'heures-CPU. C'est hors de portée de ce portable (mémoire : 15\,Go partagés avec Windows, trois alertes
mémoire cette nuit), mais adapté à la machine de calcul du projet (32\,Go) sur quelques jours, ou à un cluster.
\paragraph{Arrêt sur ce portable (23 septembre, 10\,h\,14).} À la demande de l'utilisateur, les solveurs de cubes ont été arrêtés proprement :
3\,381 cubes UNSAT au total, 50 cubes durs fermés par découpage, 1 indécis, aucun cube SAT. Par instance (UNSAT sur 325) : m00028\_sL0 325,
m16420\_sL1 324, m23221\_s0027 315, m21041\_s0025 313, m23431\_s0028 309, m06717\_s0014 307, m20723\_s0024 306, m12295\_s0015 302,
m04907\_s0007 301, m05033\_s0012 291, m29481\_s0032 288. L'énumération de la classe 29 a été arrêtée avec 31 multi-ensembles enregistrés (liste
incomplète).
\paragraph{Réénumération avec phases et comparaison.} \code{enum\_supports\_phased.py} a été relancé sur les classes de l'ancienne chaîne.
Toutes celles terminées ont exactement \textbf{un} multi-ensemble compatible, sauf la classe 0 (4) et la classe 4 (96). Comparaison avec
l'ancienne chaîne indépendante (énumération phasée de $B$ puis filtre LQ, \code{compare\_phased\_old.py}) sur 30 classes : 27 identiques ;
les 3 écarts sont attendus (classe 4 : l'ancienne liste ne contient que les 18 représentants par symétrie des 96 ; classes 25 et 32 : l'ancienne
liste contenait m08290\_sL1 et m09313\_sL1, exclus par les phases avec certificat de Farkas).

\paragraph{Pour reprendre.} Tout est reprenable : \code{cubes/LIST.txt} (file), \code{cubes/<instance>/results.txt} et \code{subcubes.txt}
(résultats), \code{sh cube\_runner\_inc.sh 5000} (un runner par cœur, environ 300\,Mo chacun), marqueur \code{FROM\_SPLIT} pour aller
directement au découpage. Certification : \code{cert\_sym.sh} (instance entière) ; pour les cubes, une CNF = base + clauses unitaires du cube.

%==================================================================
\section{Plan de calcul et reproduction (23 septembre 2026)}\label{sec:plan}
\emph{Pour la personne qui fournira le calcul et pour celle qui vérifiera. Chiffres mesurés sur les exécutions de la nuit (fichiers
\code{cubes/*/results.txt}, \code{NUIT\_JOURNAL.md}).}

\subsection{Ce que décide le calcul complet}
Chaque étape est exhaustive : graphes $F$ (26 sur 156), classes de phases (46), multi-ensembles de supports par classe, puis étage quotient
découpé en cubes qui couvrent tous les cas (chaque orbite indépendante a exactement deux doubles). Si tout va au bout :
\begin{itemize}[nosep]
\item \textbf{tout est UNSAT} : aucun $\srg(99,14,1,2)$ n'a d'automorphisme d'ordre 3 (sous réserve des points de la section suivante) ;
\item \textbf{un problème quotient est SAT} : ce n'est \emph{pas} encore une réponse, car l'étage quotient ignore les phases du bloc $C$.
Il faut alors l'étage des phases (\code{stage\_p.py}, prêt et testé sur le graphe tour) : soit on obtient un graphe, vérifiable en une seconde
($A^2+A=12I+2J$), soit ce $Q$ est exclu et il faut traiter les autres solutions quotient de l'instance. Ce cas n'est pas chiffré.
\end{itemize}
{\sloppy Indices à ce jour : 10 problèmes quotient entièrement UNSAT, 3\,381 cubes UNSAT à l'arrêt (\code{MANIFEST.txt}), \textbf{aucune solution nulle part}, aucune solution
invariante pour l'instance la plus symétrique. Cela rend UNSAT très probable, sans le démontrer.\par}

\subsection{Niveau de confiance de chaque maillon}
\begin{center}\small\begin{tabular}{p{0.28\textwidth}p{0.30\textwidth}p{0.34\textwidth}}
maillon & garanti par & ce qui manque\\\hline
absence de point fixe de $g$ & articles de revues (Makhnev--Minakova 2004, Behbahani--Lam 2011) & --\\
$\tau=6$, graphes $F$, classes & démonstrations (dont une preuve autonome de $\tau=6$), certificats de Farkas & relecture humaine\\
dépôts GitHub de tiers & -- & aucun résultat n'en dépend (contexte seulement)\\
lemmes des entrées, LQ, niveau 2.5 & démonstrations (\S\ref{sec:l25}, \S\ref{sec:nuit}) & \textbf{relecture humaine du niveau 2.5}\\
cassage de symétrie, couverture des cubes & lemme~\ref{lem:symlex}, validations exhaustives & relecture humaine\\
encodages CNF & validations sur graphes réels et témoins SAT & (preuve formelle hors de portée)\\
UNSAT de l'étage quotient & DRAT + drat-trim (9 instances) & DRAT cube par cube pour les autres\\
\textbf{complétude des énumérations} & « infaisable » de CP-SAT ; certificat DRAT (\code{cert\_enum.py}) & certifiée pour 7 classes
(0, 4, 15, 25, 30, 32, 42) ; \textbf{35 autres complètes à certifier ; 29, 43, 44, 45 incomplètes}
\end{tabular}\end{center}

\paragraph{Certifier une énumération.} L'« infaisable » final de CP-SAT n'est pas vérifiable : sur un modèle test, CP-SAT conclut dès le
pré-traitement (symétries, réductions duales) et n'émet aucune preuve, bien qu'il en expose l'option. \code{cert\_enum.py} réécrit donc le modèle
d'énumération avec phases entièrement en CNF, contrainte par contrainte : supports triés, $G_0$, TY, sommes de lignes, doubles, L1, C3, colonne
phasée par orbite (parmi celles qu'autorise le filtre de type), égalités résiduelles, condition des triplets. On y ajoute la clause bloquante de
chaque multi-ensemble trouvé ; l'insatisfiabilité, prouvée par CaDiCaL et vérifiée par drat-trim, certifie que la liste est complète. Une
contrainte de symétrie est ajoutée : deux orbites de même support ont des indices de colonne rangés dans l'ordre. Elle est valide, car échanger
deux telles orbites (bloc $C$ compris) transforme une solution en solution.
\emph{Validation dans les deux sens.} (i) Toute solution de la CNF est relue par un évaluateur Python indépendant qui teste chaque contrainte du
modèle ; (ii) chaque multi-ensemble trouvé par CP-SAT, débloqué seul, redonne SAT avec exactement ce multi-ensemble. Classe 15 (4,6\,M clauses) :
débloquée SAT (0 violation), bloquée UNSAT. Classe 0 : chacun des 4 multi-ensembles débloqué seul redonne SAT avec lui-même (0 violation), les 4
bloqués donnent UNSAT.
\emph{Certificats} (CaDiCaL 3.0.1, preuve DRAT binaire, drat-trim VERIFIED) :
\begin{center}\footnotesize\begin{tabular}{rrlrrr}
classe & multi-ens. & CNF (sha256, début) & résol. & preuve & vérif.\\\hline
0 & 4 & \code{f791a54edc40d441} & 202\,s & 27\,Mo & 174\,s\\
4 & 96 & \code{32d5e21000cfa5d5} & 120\,s & 23\,Mo & 38\,s\\
15 & 1 & \code{6e4274b135ed5b96} & 272\,s & 30\,Mo & 224\,s\\
25 & 1 & \code{4a955e8e9e58b6b9} & 1\,179\,s & 84\,Mo & 330\,s\\
30 & 1 & \code{c2e9662631bbf6cc} & 140\,s & 45\,Mo & 67\,s\\
32 & 1 & \code{4cfb3bd63aa6c23d} & 180\,s & 61\,Mo & 79\,s\\
42 & 1 & \code{b096ddbee272e204} & 863\,s & 76\,Mo & 611\,s
\end{tabular}\end{center}
Le coût est de quelques minutes à une demi-heure par classe. Réénumération avec phases : 42 classes sur 46 ont une liste complète (un seul
multi-ensemble compatible par classe, sauf les classes 0 et 27, qui en ont 4, et la classe 4, qui en a 96) ; restent incomplètes 29 (31 trouvés),
43 (2 trouvés), 44 et 45. Les 35 listes complètes non encore certifiées le seront par \code{cert\_enum\_all.sh} (image \code{c3box}), environ
6 à 10\,h sur une machine.

\subsection{Coût mesuré et estimation}
Avec la méthode incrémentale, les passes à 5\,000 puis 50\,000 conflits règlent 85 à 95\,\% des 325 cubes d'une instance ; un cube dur coûte
en moyenne 1 à 45\,min selon l'instance (moyennes mesurées : 1\,min pour m00028\_sL0, 15 pour m16420\_sL1, 22 pour m12295\_s0015, 45 pour
m23431\_s0028). m00028\_sL0 a été réglée entièrement en environ 3\,h de CPU.
\begin{center}\small\begin{tabular}{lr}
poste & estimation (heures-CPU)\\\hline
11 instances en cours (environ 200 cubes durs restants) & 50--100\\
30 instances non commencées (dont les 18 de la classe 4), 3--6\,h chacune & 90--180\\
certification DRAT de tous les UNSAT & du même ordre que la résolution\\
classes 29, 43, 44, 45 & \textbf{inconnu} (principal risque)
\end{tabular}\end{center}

\subsection{Passage à l'échelle}
\begin{itemize}[nosep]
\item Environ 40 instances $\times$ 325 cubes, soit environ 13\,000 tâches indépendantes, plus les sous-cubes (25 par nœud dur, eux aussi
indépendants). Chaque tâche : 1 cœur, environ 300\,Mo. Machine de calcul du projet (16 threads, 32\,Go) : 12 à 14 runners, un à deux jours pour la
partie chiffrée. 100 cœurs : 2 à 3\,h. 1\,000 cœurs : moins d'une heure, la durée étant alors bornée par le cube le plus long (à redécouper).
\item Limite : l'énumération CP-SAT d'une classe n'utilise bien que 8 à 16 threads ; pour des centaines de cœurs, il faut la découper elle aussi
(par exemple selon le support de la première orbite), à concevoir.
\item Les clusters n'acceptent généralement pas Docker : Apptainer/Singularity importe les archives d'images telles quelles.
\end{itemize}

\subsection{Reproduction octet pour octet}
{\sloppy
Trois niveaux, qu'il faut distinguer :
\begin{enumerate}[nosep]
\item \textbf{Identique à l'octet, vérifié par SHA-256.} Fichiers d'instances ; CNF (vérifié : m16420\_sL0 donne le même SHA-256
\code{a706468daf92ee33\ldots} avec deux graines de hachage Python différentes et dans trois générations indépendantes) ; images Docker (archives
exportées, \code{docker/IMAGES.sha256}, et identifiants d'image à retrouver après \code{docker load}, \code{docker/IMAGES.txt}) ; résultats sous
forme canonique triée (\code{manifest.py}, \code{MANIFEST.txt}). Vérification croisée : les listes de supports calculées par trois méthodes
différentes (LQ seule, réduction par $\operatorname{Aut}(F)$, phases) ont des SHA-256 canoniques identiques sur les masques 28, 8290, 9313, 13452
et 29448.
\item \textbf{Déterministe si l'on configure à l'identique.} CaDiCaL sur un thread, avec des budgets de \emph{conflits} et non de secondes, redonne
le même verdict, et même la même preuve DRAT octet pour octet. Les exécutions de cette nuit ont utilisé des limites de temps et ont été
interrompues puis reprises : leurs journaux ne se rejouent pas à l'identique, leurs verdicts si.
\item \textbf{Non reproductible par nature.} CP-SAT à plusieurs threads (ordre de découverte, durées). Seul l'\emph{ensemble} final compte, d'où
la forme canonique.
\end{enumerate}
Critère scientifique retenu, plus fort qu'un simple rejeu : \textbf{même CNF (SHA-256) et certificat DRAT vérifié par drat-trim}. Ce contrôle ne
dépend ni du solveur, ni de la machine, ni de l'ordre d'exécution.\par}

\subsection{Ordre des opérations recommandé}
{\sloppy
\begin{enumerate}[nosep]
\item Charger les images : \code{docker load < docker/c99audit.image.tar.gz} (idem \code{proofbox}, \code{texbox}) ; comparer les
identifiants avec \code{docker/IMAGES.txt}. À défaut : \code{docker/Dockerfile.c99audit}, \code{docker/proofbox/Dockerfile} (versions figées,
mais binaires pas garantis identiques).
\item Lancer les runners de cubes (un par cœur) :
\code{docker run -d --cpus 1 -m 1500m -v "\$PWD:/w" -w /w c99audit sh cube\_runner\_inc.sh 5000} ; la file est \code{cubes/LIST.txt}.
\item En parallèle, les énumérations des classes 29, 43, 44, 45 :
\code{python enum\_supports\_phased.py <classe> 3600 1000 8}, puis \code{make\_phased\_instances.py} et ajout à la file.
\item Certifier : \code{cert\_sym.sh} pour une instance entière ; pour un cube, CNF de base plus ses clauses unitaires, même vérification.
\item Certifier les énumérations : \code{sh cert\_enum\_all.sh} dans l'image \code{c3box} (génère chaque CNF, puis CaDiCaL, preuve DRAT et
drat-trim, une classe à la fois ; résultats dans \code{certs/SUMMARY\_ENUM.txt}).
\item Régénérer \code{MANIFEST.txt} et le comparer ligne à ligne avec toute autre reproduction.
\end{enumerate}\par}


%==================================================================
\appendix
\section{Tableaux (générés automatiquement depuis les fichiers de résultats)}
\subsection*{A. Les 44 graphes $F$ à Gram positif}
{\small% table A: the 44 motifs F with PSD Gram G0
\begin{longtable}{rrlll}\toprule masque & $|E|$ & degr\'es & ar\^etes $ij$ & incidence agr\'eg\'ee\\\midrule\endhead
0 & 0 & 000000 & --- & t\'emoin exact\\
16384 & 1 & 110000 & 45 & t\'emoin exact\\
96 & 2 & 211000 & 12 13 & exclu, Farkas ($y\cdot b=-6$)\\
16388 & 2 & 111100 & 03 45 & t\'emoin exact\\
28 & 3 & 311100 & 03 04 05 & t\'emoin exact\\
16420 & 3 & 111111 & 03 12 45 & t\'emoin exact\\
20496 & 3 & 221100 & 05 34 45 & exclu, Farkas ($y\cdot b=-1$)\\
4680 & 4 & 321110 & 04 13 23 34 & exclu, Farkas ($y\cdot b=-1$)\\
8290 & 4 & 222110 & 02 12 13 35 & t\'emoin exact\\
10512 & 4 & 411110 & 05 15 25 35 & exclu, Farkas ($y\cdot b=-1$)\\
24588 & 4 & 222200 & 03 04 35 45 & exclu, Farkas ($y\cdot b=-2$)\\
929 & 5 & 421111 & 01 12 14 15 23 & exclu, Farkas ($y\cdot b=-1$)\\
4649 & 5 & 222220 & 01 04 12 23 34 & exclu, Farkas ($y\cdot b=-30$)\\
4776 & 5 & 322210 & 04 12 14 23 34 & t\'emoin exact\\
9313 & 5 & 322111 & 01 12 13 24 35 & t\'emoin exact\\
12295 & 5 & 331111 & 01 02 03 34 35 & t\'emoin exact\\
26896 & 5 & 511111 & 05 15 25 35 45 & t\'emoin exact\\
29192 & 5 & 332110 & 04 23 34 35 45 & exclu, Farkas ($y\cdot b=-2$)\\
29248 & 5 & 422110 & 13 23 34 35 45 & exclu, Farkas ($y\cdot b=-1$)\\
126 & 6 & 422211 & 02 03 04 05 12 13 & exclu, Farkas ($y\cdot b=-3$)\\
1731 & 6 & 332220 & 01 02 13 14 23 24 & exclu, Farkas ($y\cdot b=-6$)\\
4653 & 6 & 332220 & 01 03 04 12 23 34 & exclu, Farkas ($y\cdot b=-2$)\\
4905 & 6 & 322221 & 01 04 12 15 23 34 & t\'emoin exact\\
8302 & 6 & 332211 & 02 03 04 12 13 35 & exclu, Farkas ($y\cdot b=-60$)\\
13452 & 6 & 432111 & 03 04 14 24 34 35 & t\'emoin exact\\
15426 & 6 & 332211 & 02 13 24 25 34 35 & t\'emoin exact\\
21041 & 6 & 222222 & 01 05 12 23 34 45 & t\'emoin exact\\
29448 & 6 & 333111 & 04 15 23 34 35 45 & t\'emoin exact\\
29706 & 6 & 422220 & 02 04 24 34 35 45 & exclu, Farkas ($y\cdot b=-60$)\\
30992 & 6 & 522111 & 05 15 25 34 35 45 & t\'emoin exact\\
4907 & 7 & 333221 & 01 02 04 12 15 23 34 & t\'emoin exact\\
5033 & 7 & 432221 & 01 04 12 14 15 23 34 & t\'emoin exact\\
12857 & 7 & 332222 & 01 04 05 12 23 34 35 & exclu, Farkas ($y\cdot b=-2$)\\
16927 & 7 & 522221 & 01 02 03 04 05 23 45 & t\'emoin exact\\
21045 & 7 & 332222 & 01 03 05 12 23 34 45 & t\'emoin exact\\
23202 & 7 & 432221 & 02 12 14 23 25 34 45 & exclu, Farkas ($y\cdot b=-3$)\\
26850 & 7 & 333221 & 02 12 13 14 25 35 45 & t\'emoin exact\\
6717 & 8 & 433222 & 01 03 04 05 12 23 25 34 & t\'emoin exact\\
15768 & 8 & 442222 & 04 05 14 15 24 25 34 35 & exclu, Farkas ($y\cdot b=-1$)\\
20723 & 8 & 433222 & 01 02 05 12 13 14 34 45 & t\'emoin exact\\
23217 & 8 & 333322 & 01 05 12 14 23 25 34 45 & exclu, Farkas ($y\cdot b=-2$)\\
29481 & 8 & 333322 & 01 04 12 15 23 34 35 45 & t\'emoin exact\\
23221 & 9 & 333333 & 01 03 05 12 14 23 25 34 45 & t\'emoin exact\\
23431 & 9 & 333333 & 01 02 03 14 15 23 25 34 45 & t\'emoin exact\\
\bottomrule\end{longtable}

}
\subsection*{B. Les 46 classes (graphe, phases) et l'énumération de $B$}
{\small% table B: the 46 (F, phase) classes up to symmetry and the enumeration of B
\begin{longtable}{rrlrrcr}\toprule classe & masque & phases $\varepsilon$ & membres & $B$ phas\'es & \'enum.\ compl\`ete & supports\\\midrule\endhead
0 & 16927 & 0000000 & 1 & 2448 & oui & 6\\
1 & 26850 & 0000001 & 4 & 20 & oui & 1\\
2 & 5033 & 0000001 & 2 & 767 & oui & 5\\
3 & 5033 & 0000000 & 1 & 250 & oui & 3\\
4 & 30992 & 000000 & 1 & 2429 & \textbf{non} & 96\\
5 & 13452 & 000000 & 1 & 2127 & \textbf{non} & 7\\
6 & 29448 & 000000 & 1 & 2005 & \textbf{non} & 10\\
7 & 15426 & 000000 & 1 & 76 & oui & 2\\
8 & 4907 & 0000001 & 2 & 8 & oui & 1\\
9 & 4907 & 0000000 & 1 & 2 & oui & 1\\
10 & 15426 & 000001 & 2 & 300 & oui & 2\\
11 & 6717 & 00000010 & 4 & 105 & oui & 2\\
12 & 4905 & 000000 & 1 & 878 & oui & 3\\
13 & 4905 & 000001 & 2 & 653 & oui & 3\\
14 & 20723 & 00000001 & 2 & 12 & oui & 1\\
15 & 21041 & 000000 & 1 & 14 & oui & 1\\
16 & 26850 & 0000011 & 2 & 16 & oui & 1\\
17 & 26850 & 0000012 & 2 & 52 & oui & 1\\
18 & 6717 & 00000100 & 2 & 67 & oui & 2\\
19 & 6717 & 00000120 & 2 & 118 & oui & 2\\
20 & 4776 & 00001 & 2 & 762 & oui & 4\\
21 & 29481 & 00000000 & 1 & 16 & oui & 1\\
22 & 4776 & 00000 & 1 & 192 & oui & 3\\
23 & 21045 & 0000001 & 4 & 32 & oui & 1\\
24 & 21045 & 0000000 & 1 & 4 & oui & 1\\
25 & 8290 & 0000 & 1 & 938 & \textbf{non} & 5\\
26 & 29481 & 00000011 & 2 & 186 & oui & 3\\
27 & 16420 & 000 & 1 & 1817 & \textbf{non} & 6\\
28 & 21045 & 0000102 & 2 & 138 & oui & 1\\
29 & 26896 & 00000 & 1 & 729 & \textbf{non} & 560\\
30 & 12295 & 00000 & 1 & 1320 & oui & 1\\
31 & 21045 & 0000101 & 2 & 137 & oui & 1\\
32 & 9313 & 00000 & 1 & 1905 & \textbf{non} & 4\\
33 & 20723 & 00000000 & 1 & 10 & oui & 1\\
34 & 29481 & 00000121 & 2 & 176 & oui & 3\\
35 & 29481 & 00000102 & 4 & 54 & oui & 3\\
36 & 23431 & 000000010 & 6 & 1 & oui & 1\\
37 & 23431 & 000000120 & 2 & 1 & oui & 1\\
38 & 23431 & 000000000 & 1 & 1 & oui & 1\\
39 & 23221 & 000000111 & 36 & 1 & oui & 1\\
40 & 23221 & 000000110 & 12 & 1 & oui & 1\\
41 & 23221 & 000001211 & 2 & 1 & oui & 1\\
42 & 28 & 000 & 1 & 637 & \textbf{non} & 1\\
43 & 16388 & 00 & 1 & 453 & \textbf{non} & 166\\
44 & 16384 & 0 & 1 & 207 & \textbf{non} & 182\\
45 & 0 & --- & 1 & 140 & \textbf{non} & 126\\
\bottomrule\end{longtable}

}
\subsection*{C. Étage quotient : les 33 problèmes issus des 35 classes complètement énumérées}\label{sec:tabC}
{\small% table C: Stage-Q instances of the 33 completely enumerated classes
\begin{longtable}{llrlrl}\toprule instance & classes & $\operatorname{rang}M$ & \'etage Q (pysat/CaDiCaL) & temps & preuve DRAT (drat-trim)\\\midrule\endhead
m04776\_s0000 & 20,22 & 6 & UNKNOWN & 300\,s & ---\\
m04776\_s0001 & 20,22 & 6 & UNKNOWN & 300\,s & ---\\
m04776\_s0002 & 20,22 & 6 & UNKNOWN & 300\,s & ---\\
m04776\_s0003 & 20 & 6 & UNKNOWN & 300\,s & ---\\
m04905\_s0004 & 12,13 & 6 & UNKNOWN & 300\,s & ---\\
m04905\_s0005 & 12,13 & 6 & UNKNOWN & 300\,s & ---\\
m04905\_s0006 & 12,13 & 6 & UNKNOWN & 300\,s & ---\\
m04907\_s0007 & 8,9 & 5 & UNKNOWN & 300\,s & ---\\
m05033\_s0008 & 2,3 & 6 & UNKNOWN & 300\,s & ---\\
m05033\_s0009 & 2,3 & 6 & UNSAT & 55\,s & \textbf{v\'erifi\'ee}\\
m05033\_s0010 & 2 & 6 & UNSAT & 45\,s & \textbf{v\'erifi\'ee}\\
m05033\_s0011 & 2 & 6 & UNSAT & 45\,s & \textbf{v\'erifi\'ee}\\
m05033\_s0012 & 2,3 & 6 & UNKNOWN & 300\,s & ---\\
m06717\_s0013 & 11,18,19 & 6 & UNSAT & 39\,s & \textbf{v\'erifi\'ee}\\
m06717\_s0014 & 11,18,19 & 6 & UNKNOWN & 300\,s & ---\\
m12295\_s0015 & 30 & 6 & UNKNOWN & 300\,s & ---\\
m15426\_s0016 & 7,10 & 6 & UNKNOWN & 300\,s & ---\\
m15426\_s0017 & 7,10 & 6 & UNKNOWN & 300\,s & ---\\
m16927\_s0018 & 0 & 5 & UNSAT & 37\,s & \textbf{v\'erifi\'ee}\\
m16927\_s0019 & 0 & 5 & UNKNOWN & 300\,s & ---\\
m16927\_s0020 & 0 & 5 & UNKNOWN & 300\,s & ---\\
m16927\_s0021 & 0 & 5 & non termin\'e & --- & ---\\
m16927\_s0022 & 0 & 5 & non termin\'e & --- & ---\\
m16927\_s0023 & 0 & 5 & UNSAT & 16\,s & \textbf{v\'erifi\'ee}\\
m20723\_s0024 & 14,33 & 5 & UNKNOWN & 300\,s & ---\\
m21041\_s0025 & 15 & 5 & non termin\'e & --- & ---\\
m21045\_s0026 & 23,24,28,31 & 6 & UNKNOWN & 300\,s & ---\\
m23221\_s0027 & 39,40,41 & 6 & non termin\'e & --- & ---\\
m23431\_s0028 & 36,37,38 & 5 & UNKNOWN & 300\,s & ---\\
m26850\_s0029 & 1,16,17 & 6 & non termin\'e & --- & ---\\
m29481\_s0030 & 26,34,35 & 6 & UNSAT & 63\,s & \textbf{v\'erifi\'ee}\\
m29481\_s0031 & 26,34,35 & 6 & UNSAT & 100\,s & \textbf{v\'erifi\'ee}\\
m29481\_s0032 & 21,26,34,35 & 6 & UNKNOWN & 300\,s & ---\\
\bottomrule\end{longtable}
}

\section{Reproduction}
Tout est dans \code{Desktop/c3\_recherche\_annexe}. Images Docker : \code{c99audit} (Python 3.12, sympy 1.14, numpy 2.5.3, python-sat 1.9.dev15,
ortools 9.15, networkx), \code{proofbox} (CaDiCaL 3.0.1 et drat-trim compilés depuis les sources officielles), \code{texbox} (LaTeX).
{\small\begin{verbatim}
python check_all.py                      # 26 temoins + 18 Farkas (stdlib seule)
python audit_verify.py                   # verification independante A/B/C (networkx)
python validate_quotient.py              # modele quotient sur srg(243,22,1,2) et tour 3x3
python validate_full_model.py            # modele complet sur srg(243,22,1,2)
python level2_phased.py 10 ; python level25.py 10 ; python reduce_classes.py
python enum_all.py 120                   # enumeration des B (46 classes)
python sweep_q.py 5 300 complete         # etage quotient (33 problemes)
python cert_q.py stageQ_runs/<nom>.in.json 0 gen     # CNF exacte + SHA-256
sh cert_check.sh <nom> 1800                         # CaDiCaL + drat-trim
python spectral_test.py                  # condition spectrale (resultat negatif)
\end{verbatim}}
L'inventaire des fichiers avec empreintes SHA-256 est dans \code{rapport/INVENTAIRE.txt}.

\section{Sources}
\begin{sloppypar}
\begin{itemize}[nosep]
\item A.\,A.~Makhnev, I.\,M.~Minakova, Diskretn.\ Mat.\ 16:1 (2004) 95--104. M.~Behbahani, thèse, Concordia, 2009 ; M.~Behbahani, C.~Lam, Discrete Math.\ 311 (2011).
\item Y.~Ishida, \emph{No involutions in the missing Moore graph}, arXiv:2606.29183v2 (2026), prop.~4.4 et 8.7.
\item \code{infinityscroll/conway99-order3-f27} (14 août 2026) ; \code{ibenarb/conway99-research} ; \code{Lundle22/conway-99-cyclic-lifts} (GitHub).
\item Dossier de départ \code{codex\_c3\_demarrage} (ChatGPT) : lemme de l'hexagone et généralisation, revérifiés ici.
\item M.~Codish, A.~Miller, P.~Prosser, P.\,J.~Stuckey, \emph{Breaking symmetries in graph representation}, IJCAI 2013.
\item E.\,R.~Berlekamp, J.\,H.~van Lint, J.\,J.~Seidel (1973) : $\srg(243,22,1,2)$, utilisé comme contrôle.
\item M.\,J.\,H.~Heule, \code{drat-trim} ; A.~Biere et al., CaDiCaL.
\end{itemize}
\end{sloppypar}
\end{document}
